\documentclass[11pt]{article}
\usepackage{../common/pbstyle}
%% Common header for TD / TP / project sheets (same layout as the other sheets of the course)
\newcommand{\sheethead}[2]{%
\begin{center}
{\LARGE\bfseries Ensemble Methods}\\[0.2cm]
{\Large #1}\\[0.3cm]
{\large M2 Computer Science -- MALIA}\\[0.4cm]
Guillaume Metzler\\
Institut de Communication (ICOM)\\
Universit\'e de Lyon, Universit\'e Lumi\`ere Lyon 2\\
Laboratoire ERIC UR 3083, Lyon, France\\[0.2cm]
\texttt{guillaume.metzler@univ-lyon2.fr}
\end{center}
\vspace{0.3cm}
\begin{center}\begin{minipage}{0.88\linewidth}\small\textbf{Abstract.} #2\end{minipage}\end{center}
\vspace{0.3cm}}

\renewcommand{\course}{Ensemble Methods -- Project PAC-Bayes}
\usepackage{needspace}
\newcommand{\code}[1]{\texttt{#1}}

\begin{document}
\sheethead{Project -- PAC-Bayesian Learning of Weighted Majority Votes}{The goal of this project is
to study, implement and compare learning algorithms for weighted majority votes that come with
PAC-Bayesian risk certificates. You will build a small library, run a rigorous experimental study
on the datasets of the course, and explore one research direction based on a recent paper. The
project mobilizes Sections~3 to~7 of the lecture notes (in particular the self-bounding part, Sections~6--7) and notebooks~3 to~5.}

Groups of two persons are allowed to do this project. The work shall be submitted at the following
mail address
\begin{center}\texttt{guillaume.metzler@univ-lyon2.fr}\end{center}
before the deadline mentioned in your program. The submission shall contain (i)~the Python code
(package or notebooks) used to produce \emph{all} the results of the report, (ii)~a report in PDF
(10 to 15 pages, figures and tables included) presenting the methods, the experimental protocol, the
results and their analysis. The code shall reproduce the results presented in the report.

%% ------------------------------------------------------------------
\section{Datasets}

You will use the datasets of the course (\code{datasets.zip}, loaded with \code{prepdata.py}), with
\textbf{at least 15 datasets}. As in the first project, they can be divided into two groups
(balanced / imbalanced), which may require different performance measures. In addition, you will use
\code{make\_moons} and \code{make\_classification} from \code{scikit-learn} for illustrations.

%% ------------------------------------------------------------------
\section{Part A -- A library for PAC-Bayesian majority votes (mandatory)}

Starting from \code{pbtools.py} (you may rewrite it entirely), implement a Python class
\code{PBMajorityVote} with a \code{scikit-learn}-like interface (\code{fit}, \code{predict},
\code{certificate}) which:
\begin{enumerate}
\item builds a set of $n$ voters by bagging (decision trees, with a parameter for the size of the bootstrap
      samples and for the randomization of the trees) and keeps track of the out-of-bag samples;
\item computes the out-of-bag statistics: individual losses, tandem losses and disagreements, with the sizes
      $n_1$ and $n_2$ of the smallest OOB sets and intersections;
\item learns a posterior $\rho$ among: uniform, minimizer of the first-order PAC-Bayes-$\lambda$ bound
      \citep{thiemann2017strongly}, minimizer of the tandem bound \citep{masegosa2020second}, minimizer of a
      PAC-Bayesian C-bound \citep{germain2015risk,viallard2021self};
\item returns the certificates (first-order, tandem and C-bound) of the learned vote, computed with the kl
      inversion, for a confidence $\delta$ given by the user.
\end{enumerate}
The code must be tested: check your kl inversion against a brute-force search, and check the identity
$R(G_\rho)=e_\rho+\frac12d_\rho$ on a test sample. Explain in the report how the out-of-bag construction
makes the certificates valid (Theorem~7.1 of the lecture notes).

%% ------------------------------------------------------------------
\section{Part B -- Experimental study (mandatory)}

\paragraph{Protocol.} For each dataset, use at least 5 random train/test splits (stratified). Hyperparameters
that are tuned (depth of the trees, number of features per split, \dots) must be tuned by cross-validation on
the training set only. Report means and standard deviations, and use a statistical test (e.g.\ Wilcoxon signed-rank
test over the datasets) when you compare two methods.

\paragraph{Methods to compare.}
\begin{itemize}
\item a random forest (\code{scikit-learn}) and AdaBoost, as baselines without certificates;
\item your \code{PBMajorityVote} with the four posteriors (uniform, FO, TND, C-bound);
\item a single decision tree and a linear SVM, to measure the benefit of the vote.
\end{itemize}

\paragraph{Expected results and questions.}
\begin{enumerate}
\item A table of the test risks (and of a balanced measure for the imbalanced group) and a table of the
      certificates. Which method gives the best vote? the best certificate? Are the certificates non-vacuous
      (smaller than $0.5$) on all datasets?
\item The influence of the number of voters $n\in\{10,50,100,200,500\}$ and of the bootstrap size on the test
      risk and on the certificates (use figures).
\item The influence of the diversity: for a few datasets, plot the Gibbs risk, the disagreement, the risk of the
      vote and the bounds as functions of the number of features tried at each split. Discuss the
      strength/diversity trade-off in the light of the C-bound.
\item The computation times of the different methods.
\end{enumerate}

%% ------------------------------------------------------------------
\section{Part C -- A research direction (choose one)}

\begin{description}
\item[C1 -- Boosting with certificates.] Learn the voters with AdaBoost or gradient boosting on one part of the
      data and learn their weights by minimizing the tandem bound or the C-bound on the other part. Compare with the
      weights given by the boosting algorithm, in terms of risk, certificate and margins \citep{schapire1998boosting,roy2016column}.
\item[C2 -- Kernels as priors.] Following \citet{letarte2019pseudo} (and the first project of the course), use random
      Fourier features as voters with the kernel as prior. Learn a pseudo-posterior over the features (Gibbs posterior,
      or minimization of a second-order bound), and compare with an RBF SVM and with the gradient boosting with random
      Fourier features of the first project.
\item[C3 -- Tighter bounds for the vote.] Read \citet{wu2021chebyshev} and implement the Chebyshev--Cantelli
      bound with the PAC-Bayes-Bennett inequality (CCPBB). Compare it with the tandem bound and the C-bound, as
      certificates and as learning objectives.
\item[C4 -- Linear classifiers and data-dependent priors.] Implement PBGD \citep{germain2009pac} with Gaussian
      posteriors, with a prior centred at $\mathbf 0$ and with a data-dependent prior
      \citep{ambroladze2006tighter,parrado2012pac}. Compare the certificates and the risks with those of a linear SVM,
      as a function of the proportion of data used to learn the prior.
\item[C5 -- Stochastic majority votes.] Read \citet{zantedeschi2021learning} and implement their Dirichlet posterior
      over the weights of the vote on a small number of voters (e.g.\ 20 trees). Compare the certificates with those of Part~B.
\end{description}
For the chosen direction, the report must contain a short presentation of the method (with the key equations
and their justification), your implementation choices, and an experimental comparison on at least 5 datasets.

%% ------------------------------------------------------------------
\section{Evaluation}

\begin{center}
\begin{tabular}{lr}
\toprule
Part A: correctness and quality of the code, tests & 5 points \\
Part B: protocol, results, figures and analysis & 8 points \\
Part C: understanding of the method, implementation, experiments & 5 points \\
Quality of the report (clarity, notation, references) & 2 points \\
\bottomrule
\end{tabular}
\end{center}

\section*{Suggestions for the report}

\paragraph{Introduction.} Present the context (ensemble methods, need for guarantees) and the objectives.

\paragraph{Methodology.} Introduce the notation, then the PAC-Bayesian bounds you use and the learning
algorithms, with the key equations: explain what is minimized and why the certificate is valid. Present the
performance measures and the experimental protocol.

\paragraph{Results.} Tables and figures, always with a sentence explaining what they show. Compare methods that
are comparable, and discuss both the risk \emph{and} the tightness of the certificates.

\paragraph{Conclusion.} Summarize your findings, their limits and possible extensions.

\bibliographystyle{plainnat}
\bibliography{../common/pacbayes}
\end{document}
